\documentclass[11pt,a4paper]{article}

% ---------------------------------------------------------------
%  Cover Letter -- Stanisław Szufa
%  EPFL IC Faculty Search 2026
% ---------------------------------------------------------------
\usepackage[T1]{fontenc}
\usepackage[utf8]{inputenc}
\usepackage{mathpazo}
\renewcommand{\ttdefault}{lmtt}  % outline typewriter font (avoids bitmap ectt)
\usepackage[hmargin=2.5cm,vmargin=1.9cm]{geometry}
\usepackage{microtype}
\usepackage{xcolor}
\usepackage[hidelinks]{hyperref}

\definecolor{accent}{RGB}{160,20,30}   % same accent as the other documents
\setlength{\parindent}{0pt}
\setlength{\parskip}{5pt}
\pagestyle{empty}

\begin{document}

\begin{center}
  {\LARGE\bfseries\color{accent} Stanisław Szufa}\\[4pt]
  {\small
  University of Geneva, Switzerland \;$\cdot$\;
  \href{mailto:stanislaw.szufa@unige.ch}{stanislaw.szufa@unige.ch} \;$\cdot$\;
  +41 76 472 16 09}
\end{center}
{\color{accent!40}\rule{\linewidth}{0.5pt}}

\begin{minipage}[t]{0.6\linewidth}
Prof.\ Nate Foster\\
Chair of the Faculty Recruiting Committee\\
School of Computer and Communication Sciences\\
EPFL, Lausanne, Switzerland
\end{minipage}%
\hfill
\begin{minipage}[t]{0.35\linewidth}\raggedleft
19 October 2026
\end{minipage}

\vspace{2pt}
Dear Professor Foster and members of the Committee,

I am applying for the position of Tenure Track Assistant Professor in the
School of Computer and Communication Sciences. I work in computational social
choice, which studies voting and collective decisions with the tools of
computer science, and I build the methods, data, and software needed to test
whether a collective decision represents the people it affects. I am a
postdoctoral researcher at the University of Geneva; before that, I was a CNRS
postdoc at LAMSADE, Paris Dauphine-PSL, with Jérôme Lang, and I did my PhD
with Piotr Faliszewski at Jagiellonian University.

My best-known contribution is the \emph{map of elections}. I proposed a
distance under which elections can be compared regardless of how their
candidates and voters are labeled, and used it to place real and synthetic
elections in a common space, so that one can see which models of synthetic
votes resemble real data. The framework (\emph{Artificial Intelligence},
2025) now extends to stable matching, tournaments, and fair division. I also initiated Pabulib, the largest open library of
participatory budgeting data, and I am part of the team that convinced Polish
cities to adopt the Method of Equal Shares, a proportional rule for
participatory budgets. My thesis received the Polish Prime Minister's Award
(2025); I also hold a START fellowship of the Foundation for Polish Science
and the AAMAS-2023 Best Student Paper Award. I have published 5 journal
articles and 28 full conference papers, mostly at IJCAI, AAMAS, AAAI, and
NeurIPS.

At EPFL, I will extend these tools to collective decisions in which opinions
are expressed in text and aggregated by AI, such as summaries of public
consultations and consensus statements. My agenda has three connected directions: a geometry of decision
problems that covers text as well as ballots; participatory democracy as a
living laboratory, with a benchmark of votes over texts built from
participatory budgets and the UN General Assembly record; and principled
methods for auditing whether AI-written outcomes represent the people they
summarize. I plan to fund the group through the SNSF and an ERC Starting Grant.

IC does not yet have a group in computational social choice, but my agenda draws on its
strengths: sublinear
algorithms for computing distances on very large elections (Michael Kapralov),
proof of personhood and secure e-voting for digital participatory processes
(Bryan Ford), and NLP and open multilingual language models such as Apertus
for auditing representation across languages (Antoine Bosselut).
Swiss direct democracy, and Aarau's use of the Method of Equal Shares, make
Switzerland a natural setting for field studies with municipalities.

In teaching, I am prepared to take on core
bachelor courses such as algorithms, discrete mathematics, and theory of
computation, as well as game theory at the master level, and I would like to
develop a master course on algorithms for collective decisions and a seminar
on AI and democracy. I have co-supervised more than ten bachelor's and
master's students, several of whom became coauthors of papers at AAAI,
AAMAS, ECAI, and NeurIPS.

Enclosed are my CV with references, publication list, research statement,
and teaching statement. Thank you for considering my application.

Sincerely,\\[14pt]
Stanisław Szufa

\end{document}
